\section{讨论}\label{sec:6}

前述各节提出的形式模型与实现引入了一种面向动态可组合性的编程范式。本节考察该范式如何延伸至更广泛的工程关切，并讨论其中的设计张力与开放问题。

\subsection{系统边界}\label{sec:6.1}

第\ref{subsec:3.1}节中的每个效应都携带一个逆，而该逆的含义由系统边界决定。边界将系统运行所处环境分为两部分。(1) 当系统能够独占修改某个位置并恢复修改前的状态时，该位置位于边界内部，因此对其的操作被记录在 $\Gamma$ 中且可后续恢复。(2) 当上述任一能力不满足时，该位置位于边界外部，因此对其的操作表现为 $\mathrm{id}_\Gamma$，既不被记录也不被恢复。本节展开该边界的性质及其对恢复的后果。

\subsubsection{来自协效应的边界}

协效应通过将外部位置具体化来移动边界：它将对某位置的所有访问限制为该位置提供的一组操作，并为每个操作提供逆，因此原本表现为 $\mathrm{id}_\Gamma$ 的操作得以记录在 $\Gamma$ 中并被恢复。因此边界是按位置而非按媒介划定的，因为前述两种能力都是位置的属性，且具体化改变的是位置的访问方式而保持其媒介不变。例如，当仅由系统写入某内存区域时，该区域位于边界内部；当其他进程也写入时，则位于边界外部。当仅系统能访问某文件时（如私有路径下的临时文件），该文件位于边界内部；当该路径可被其他程序读写时，则位于边界外部。移动边界本身是一种权衡：环境是否为某位置提供可逆语义，以及每次访问时提供这些语义的代价有多大。我们将在第\ref{sec:6.7}节讨论这所暗示的协同设计。

\subsubsection{获取与发射}

穿越边界的操作通常分为两个阶段。(1) 在获取阶段，操作获得访问权并在边界内部安装一条记录：\texttt{open} 安装一个由 \texttt{close} 移除的描述符，\texttt{malloc} 预留一个由 \texttt{free} 释放的内存块，\texttt{fork} 启动一个由 \texttt{kill} 终止的子进程。该记录本身是将该位置具体化的协效应的一部分，例如其维护的映射中的一条条目，而安装该条目是一个可逆效应。该记录同时也是数据可以离开的通道。(2) 在发射阶段，操作通过该通道推送数据，例如 \texttt{write} 将字节交予文件或 \texttt{send} 将数据报送上线路，该推送表现为 $\mathrm{id}_\Gamma$，将数据留在其他方可以读写的地方。因此两个阶段落在边界的两侧：获取留在边界内部，而发射穿越到外部。

\subsubsection{扣留与补偿}

必须从发射中恢复的系统有两种可用方案。一种是在产生该发射的状态确定将被保留之前扣留该发射，这是回滚恢复中的输出提交问题\cite{ref48}。另一种是补偿\cite{ref49}：一种将状态恢复到应用所提供的等价关系的动作，该等价关系粗于定义\ref{def:3.33}中的 $\simeq$，例如删除已创建的文件或退还已收取的费用。此类动作以与逆相同的后进先出顺序组合，因此第\ref{subsec:3.1}节的组合性适用于它们。但元理论不适用：定义\ref{def:4.13}的交换性是针对 $\simeq$ 证明的，必须针对更粗的等价关系重新建立。

\subsection{服务多路复用}\label{sec:6.2}

OSGi\cite{ref50}等动态组件平台围绕服务组织组合：服务是提供者在某个接口下发布、消费者绑定到的功能单元。Cordis 协效应模型呼应了这一概念，服务对应于某个键背后的接口。提供服务的组件是其提供者，注入服务的组件是其消费者。单个服务可由多个提供者实现，这种多重性可以两种形式实现。(1) 独占绑定：多个实现共享一个接口，但同一时刻至多绑定一个；编排器选择绑定哪个实现，在它们之间切换需要卸载一个提供者并加载另一个，瞬时扰动每个消费者的依赖。(2) 服务代理：作为该接口入口点的中心服务被后端提供者和消费者共同注入，因此多个提供者共存，代理在其中分派每个请求。与独占绑定相比，代理吸收了这种扰动：更新后端提供者时代理保持不变，因此消费者看不到其依赖的变化，也不会触发重新加载。

服务代理支撑三种能力：负载均衡、滚动更新和跨进程调用。

\subsubsection{负载均衡}

当多个提供者共存时，代理根据可配置的策略（例如轮询、最轻负载、延迟加权）或消费者指定的显式目标在其间分派请求。由于提供者是普通组件，可以添加或移除以伸缩容量；每个提供者通过可逆效应向代理注册，因此卸载它会恢复注册并自动从代理的路由集中移除。

\subsubsection{滚动更新}

在运行时升级服务实现归结为受控的提供者转换\cite{ref51,ref52}。执行转换时，新提供者作为额外纤程被加载并向代理注册；一旦其变为 ACTIVE，流量逐渐从旧提供者转移到新提供者（例如通过调整选择权重），旧提供者在不再承载进行中的请求后被卸载。这种提供者转换将传统上是基础设施级操作（例如容器编排、蓝绿部署）的事务变为应用级组合模式。

\subsubsection{跨进程调用}

服务代理还可以跨进程边界应用\cite{ref53}。每个进程托管自己的 Cordis 上下文及本地提供者；一个协调组件将它们连接起来，将每个视为远程提供者。跨进程服务访问由保留接口的 RPC 机制中介，使分发对消费者透明。一个注意事项是跨进程调用会产生延迟且可能在途中失败，因此以同步方式暴露会阻塞调用方。因此，旨在跨进程暴露的接口必须按照异步契约进行设计。

\subsection{访问控制与沙箱}\label{sec:6.3}

给定由独立组件组装的应用，保障应用安全需要两种互补机制：(1) 约束组件可访问的依赖，(2) 将不受信任的代码与宿主环境沙箱隔离。Cordis 通过依赖声明和拦截支持前者；后者需要外部沙箱。

\subsubsection{基于能力的访问控制}

依赖访问机制（第\ref{sec:5.1.4}节）已经构成了一种针对代理中介属性的访问控制：组件只能访问其声明的依赖；未声明的访问会引发错误。这在结构上类似于基于能力的安全\cite{ref54,ref55,ref56}，其中权限由持有引用授予而非由环境权限授予。\texttt{inject} 声明充当能力请求，上下文代理充当能力中介。由于这些请求是静态声明的，组件所需的代理中介能力的完整集合在其运行之前即已知，使编排器能在加载时审查和批准它们，而非在访问发生时才发现。

通过拦截机制，这种中介可以泛化到细粒度策略。访问控制元数据可以由上下文携带或由组件声明（定义\ref{def:3.30}），提供者在依赖被调用时查阅该元数据以决定请求是否被允许。例如，文件系统依赖可以携带声明组件可读写哪些路径的元数据，提供者根据元数据检查每次调用。由于此拦截位于上下文而非任一方的代码中，编排器可以调整它以约束任何组件对某依赖的访问而无需修改提供者，例如授予社区组件只读数据库访问权而核心组件保留完全访问权。此外，由于拦截只影响依赖的调用方式而非是否满足，它可以在运行时安装、重新配置或移除而不触发任何重新加载或扰动依赖图。

\subsubsection{沙箱化不受信任的组件}

当组件的代码不可信任时，语言级访问控制是不够的，因为拥有宿主运行时访问权的恶意组件可以直接到达底层对象，使此类检查失效。沙箱需要一种超出语言级手段范围的执行边界，例如软件故障隔离\cite{ref57}、独立语言运行时、沙箱进程或虚拟化容器\cite{ref58}。无论采用何种机制，不受信任的组件在其自己的沙箱上下文中运行，通过桥接访问宿主提供的依赖，泛化了第\ref{sec:6.2}节的跨进程调用：相同的透明性论证使得这种桥接访问与本地注入不可区分。在宿主侧，桥接是普通纤程，其能力可以通过上述访问控制来削弱。

\subsection{语言无关性与选择}\label{sec:6.4}

尽管 Cordis 以 TypeScript 实现，上下文范式是语言无关的：时空可组合性仅由其两个可组合性维度定义，因此可以在任何沿两个维度满足特定要求的语言中实现。我们依次沿每个维度分析这些要求。

\subsubsection{时间可组合性}

在最基本的层面上，时间可组合性需要闭包：可逆效应将动作与逆配对，该逆必须作为值连同其恢复的状态一起被捕获，以便在拆卸时重放。除此之外，组件的代码及加载它的副作用必须可在运行时引入和撤回。

语言如何满足第二个要求取决于其执行模型。在托管运行时中，这采取可编程模块注册表的形式，其中已加载的模块可以从注册表中驱逐并在不再被引用时被垃圾回收；例如，Node.js 就暴露了这样的注册表。\footnote{CommonJS 通过 \texttt{require.cache} 暴露模块缓存；ES 模块不提供公共驱逐 API，但模块仍可通过引擎内部接口进行管理。}原生代码不暴露模块注册表，因此引入和撤回采取显式动态链接和取消链接的形式（例如 Unix 上的 \texttt{dlopen/dlclose}，Windows 上的 \texttt{LoadLibrary/FreeLibrary}）\cite{ref59}，即将目标代码加载到运行中的进程并随后分离它。WebAssembly 根据其嵌入器选择其中一条路径：模块实例在托管嵌入器（例如 JavaScript 宿主）下由宿主的回收器回收，或在原生嵌入器（例如 Wasmtime）丢弃它时释放。跨这些机制，可逆效应模型将加载视为对上下文的效应，其逆撤销模块引入的符号、类型或处理程序的注册。

\subsubsection{空间可组合性}

空间可组合性需要组件声明其依赖的机制以及运行时提供和注入这些依赖的机制。这归结为依赖注入（DI）问题\cite{ref38}，它在两个层面表现，这两个层面在不同语言中有所不同：依赖如何被类型化以及其访问如何被中介。

在类型层面，语言应提供让开发者表达良类型依赖访问的方式。消费者通过从上下文中读取其键来获得协效应，因此上下文类型（第\ref{subsubsec:3.2.1}节）必须记录每个键的协效应。类型类（Haskell）\cite{ref60}和特征（Rust）\cite{ref61}通过让提供者从自己的模块通过实例或 \texttt{impl} 扩展上下文类型来实现这一点\cite{ref62}。TypeScript 的模块增强\cite{ref63}同样允许提供者模块将声明合并到上下文类型中。

在运行时层面，依赖访问必须被动态中介：键背后的协效应可能随提供者的加载和卸载而改变，且在不同上下文中可能被不同地解析。因此语言需要一种方式来透明地介入访问，使消费者的代码不变，例如通过 JavaScript 的 \texttt{Proxy} 对象\cite{ref64}或 Python 的描述符协议（\texttt{\_\_get\_\_}）\cite{ref65}。在没有此类原语的情况下，运行时反射\cite{ref66,ref67}可以动态中介访问，但以类型安全和开发者体验为代价。

跨两个层面，元编程设施同时提供类型化和中介。注解\cite{ref68}和装饰器将元数据附加到声明上，处理器将其展开为中介访问的访问器；编译时元编程（例如 Rust 过程宏、Scala 宏\cite{ref69}、Zig comptime）为每个依赖生成一个带类型的声明连同这样的访问器，从而免去了通用拦截原语。

\subsection{相互依赖与组件粒度}\label{sec:6.5}

在反应式协效应模型中，依赖环只会使涉及的组件永久不活跃：给定两个组件 $A$ 和 $B$，如果 $A$ 需要 $B$ 提供的键且 $B$ 需要 $A$ 提供的键，则两者的满足谓词永远不会变为真。与并发系统中依赖于调度且必须即时检测的死锁不同，此条件仅从依赖声明即可预测，因此运行时可以在组件加载时报告它。

在实践中，大多数表面上的相互依赖可以分解为更细粒度的组件以消除环。考虑两个组件：一个服务器（提供网络接口）和一个访问控制器（执行授权策略）。这两个组件双向交互：访问控制器中介到达服务器的请求，服务器暴露用于修改访问控制策略的端点。单体设计会使每个组件依赖另一个。然而，两个交互方向是逻辑上独立的关注点。分解它们产生四个组件：\texttt{server-core}、\texttt{access-control-core}、\texttt{request-mediation}（依赖两个核心以对传入请求应用访问控制）和 \texttt{policy-management}（依赖两个核心以通过服务器暴露策略修改）。通过这种方法，环被消除，因为两个核心互不依赖；只有集成组件依赖两者。

这种分解原则上总是可行的，因为每个双向交互都可以分解为独立的单向绑定，但它增加了组件的数量：在一般情况下，给定 $n$ 个相互交互的组件，集成组件的数量可以随 $n$ 二次增长，因为每对交互组件可能需要为每个交互方向一个独立组件。这不影响正确性或运行时性能（组件是轻量的），更细的粒度甚至可能是有益的：用户获得只加载其所需特定集成绑定的能力，有效增加了系统的可组合性。然而，它可能影响开发者体验：更多组件需要更多配置、更多命名以及理解依赖图的更多认知开销。

缓解这种粒度成本是工程关切而非理论关切。实用策略包括包捆绑（即将相关细粒度组件分组为单个可安装单元）、基于约定的接线（即自动连接名称或类型匹配某模式的组件）以及脚手架工具（即从声明式规范生成样板集成组件）。这些策略保留了无环模型的形式保证，同时将编写负担降低到更接近单体设计的程度。

\subsection{依赖类型化与版本化}\label{sec:6.6}

在形式模型中，依赖链接纯粹通过键同一性建立：提供键 $k$ 的组件满足任何在其依赖集中声明 $k$ 的组件。类型族 $\mathcal{V}_k$ 确保单个编译单元内的类型级一致性，但当组件独立开发和构建时，这一保证就会失效，而这在组件生态系统中是常见场景。这种失效导致两个不同的问题。

\subsubsection{接口漂移}

提供者可能在版本间修改与 $k$ 关联的接口（添加字段、更改方法签名、改变行为契约），而针对早期接口编译的消费者继续声明相同的键 $k$。该依赖在协效应层面被满足（$k \in \mathrm{dom}(\sigma)$），但运行时值不再符合消费者的期望，导致类型错误、方法未找到故障或静默的行为分歧\cite{ref70}。

\subsubsection{键冲突}

两个独立开发的提供者可能使用相同的键名 $k$ 表示完全不相关的接口。由于仅凭键同一性建立链接，期望某个提供者接口的消费者将接受另一提供者的值而不进行任何兼容性检查。与提供者和消费者至少共享共同渊源的接口漂移不同，键冲突中期望类型与实际类型之间没有任何关系，使由此产生的故障不可预测且难以诊断。

两个问题都指向同一个缺口：协效应模型仅提供名义链接（按键名），而不提供版本化或结构链接（按接口兼容性）\cite{ref71}。我们讨论三种填补该缺口的方法，从最耦合基础设施到最语言无关。

\subsubsection{键命名空间}

将键空间从 $K$ 扩展到 $K \times P$，其中 $P$ 标识定义接口的包，构造性地消除了键冲突：具有相同本地名的独立开发接口占据不同的键。这是最直接的解决方案，但也是耦合最深的：它将包命名空间嵌入形式模型本身，使系统依赖外部包注册表来确定键同一性。

\subsubsection{对等依赖}

更轻的耦合是通过宿主语言包管理器声明版本约束\cite{ref72}。这是 Cordis 当前采用的方案。组件依赖在语义上是对等依赖：组件不在内部捆绑其依赖，而是期望运行时上下文提供它们。支持对等依赖的包管理器（例如 npm）可以强制版本兼容性：如果提供某个键的包的版本落在消费者声明的对等范围之外，不兼容性在安装时被捕获而非表现为运行时故障。然而，此方案有两个局限：(1) 它依赖于提供者忠实地遵循语义版本化，这是一种不可强制执行的惯例；(2) 包管理器通常将每个依赖解析为单一版本，这阻止了在同一个应用中加载同一包的多个版本的组件。

\subsubsection{结构兼容性}

完全语言无关的方法将成员检查 $k \in \mathrm{dom}(\sigma)$ 替换为验证提供者实际接口在结构上包含消费者期望的兼容性谓词。这类似于结构子类型化\cite{ref73}：如果提供的接口是所需接口的子类型，则提供者满足消费者。挑战在于以语言无关的方式定义此谓词：结构兼容性对记录类型（宽度子类型化）是直接的，但对行为契约（例如前后条件\cite{ref74}、效应规范\cite{ref22}）变得复杂，且一旦参数多态引入有界量化则变得不可判定\cite{ref75}。

这三种方法分别解决问题的不同方面。设计一个统一依赖模型以结合这些方法，同时保留协效应模型的动态组合保证，仍然是一个开放问题。

\subsection{与语言和操作系统的协同设计}\label{sec:6.7}

第\ref{sec:6.4}节确定了宿主语言必须为上下文范式提供的最低要求。本节讨论相反的问题：与该范式协同设计的语言或操作系统在此最低要求之外还能提供什么。

\subsubsection{与语言的协同设计}

围绕上下文范式设计的语言可以在两个方面改进库实现：它赋予上下文的语义，以及它赋予效应和协效应的原语。

这样的语言可以在保留第\ref{subsec:3.3}节的上下文语义的同时使上下文再次隐式化。命令式语言已经在隐式上下文上运行每条语句，而该单一上下文既不追踪效应也不解析协效应。上下文范式转而区分多个上下文，其中操作要么修改其运行的上下文，要么从中派生出另一个上下文（定义\ref{def:3.27}）。原地实现修改环境上下文，正如命令式语言所做。派生实现则引入一个独立的上下文，语言必须为此提供构造。使上下文隐式化带来人机工程和安全两方面的益处。(1) 在库实现中，每个涉及效应或协效应的函数将上下文作为普通参数或接收者，如第\ref{sec:5.1}节所示。当语言隐式提供上下文时，函数不再需要接受它。(2) 每个上下文携带自己的生命周期状态和已提交视图（第\ref{sec:4.1}节）。库实现将上下文作为普通变量传递，因此组件可能通过闭包或全局变量误达另一组件的上下文。在那里安装的效应随后泄漏到其自身的生命周期之外，在那里读取的协效应逃脱其依赖规约。使上下文隐式化关闭了这两个漏洞。

这样的语言还可以使其编译器感知效应和协效应。(1) 对于效应，效应迭代器（定义\ref{def:4.9}）在每一步分配一个闭包以持有逆连同其恢复的状态。有了执行效应的语法，编译器可以为整个迭代发出单个状态机并将这些逆保存在其帧中。(2) 对于协效应，协效应规约可以被纳入类型系统，有两个益处。首先，依赖环在编译时被报告而非留给运行时（第\ref{sec:6.5}节）。其次，依赖可以通过其类型的结构而非仅凭键同一性进行比较，如行类型所做\cite{ref28}，这是对第\ref{sec:6.6}节结构兼容性的类型级支持。

\subsubsection{与操作系统的协同设计}

第\ref{sec:1.2.3}节观察到动态可组合性的一种粗粒度替代方案，其中操作系统以进程粒度提供时间可组合性，其上的容器编排器以服务粒度提供空间可组合性。与该范式协同设计的操作系统将支持细粒度组合，方式是将组件声明的协效应规约作为其可达范围的全部，并将其自身资源作为协效应提供。

这样的操作系统可以提供第\ref{sec:6.6}节推迟到语言外部机制的沙箱。它通过将组件限定为其声明的依赖来实现这一点，在组件加载时提供这些依赖并使组件内部无法触及其他任何东西，正如 WebAssembly 模块在实例化时从其嵌入器接收其导入\cite{ref76}。